\chapter{零测集与不可测集}
\label{chap:null-nonmeasurable-sets}

“测度为零”表示集合可由总体积任意小的可数长方体族覆盖，并不表示集合
有限、至多可数或不存在复杂结构。另一方面，外测度定义在所有子集上，却不能
在所有子集上保持平移不变和可数可加。本章分别考察这两种现象。

\section{零测集与几乎处处}

\begin{definition}
若 $N$ 可测且 $m(N)=0$，称 $N$ 为\term{零测集}。一个命题在除去某个
零测集之外的所有点成立时，称它\term{几乎处处成立}，简写为 a.e.
\end{definition}

由于 Lebesgue 测度完备，“外测度为零”已经保证可测。可数个零测集的并
仍零测；有限测度集减去零测集不改变测度。若两个函数只在零测集上不同，
测度论中的积分和多数收敛概念会把它们视为同一个对象，但点值连续性未必
保留。

零测集可以稠密，例如 $\Q$；也可以不可数、闭且没有孤立点。Cantor 集是
后一情形的标准模型。

\section{Cantor 集}

从 $C_0=[0,1]$ 开始，每一步从每个闭区间中删去开中间三分之一，令剩余
集合为 $C_n$，并定义
\[
 C=\bigcap_{n=0}^{\infty}C_n.
\]
$C_n$ 是 $2^n$ 个长度 $3^{-n}$ 的闭区间之并，故
$m(C_n)=(2/3)^n$。由下降连续性，
\[
 m(C)=\lim_nm(C_n)=0.
\]

每个 $x\in C$ 都有一个只含数字 $0,2$ 的三进制表示。第 $n$ 位取 $0$ 或
$2$，分别记录 $x$ 落在第 $n$ 次保留区间的左三分之一或右三分之一。三进制
端点虽然有两种通常表示，但小数部分的两种表示中至多一种只使用 $0,2$。
这里允许 $1=0.222\ldots{}_3$，且不一律取尾部为2的表示：例如
$2/3=0.2000\ldots{}_3$，而它的另一表示 $0.1222\ldots{}_3$ 含有1。
嵌套区间的左右地址消除了这种选择上的歧义。因此映射
\[
 (\varepsilon_1,\varepsilon_2,\ldots)\in\{0,1\}^{\Npos}
 \longmapsto \sum_{n=1}^{\infty}\frac{2\varepsilon_n}{3^n}
\]
是二元序列集到 $C$ 的双射。特别地，$C$ 不可数。

其余拓扑性质也可从构造直接读出。集合 $C$ 是闭集列 $C_n$ 的交，因而闭。
若非空开区间 $I$ 包含于 $C$，它也包含于每个 $C_n$。区间不能跨过
$C_n$ 的分支之间的空隙，故必须落在一个长度为 $3^{-n}$ 的分支内。
令 $n\to\infty$ 与 $I$ 的正长度矛盾，所以 $C$ 无内点。另一方面，给定 $x\in C$ 和
$n$，改变其三进制表示中某个晚于第 $n$ 位的数字，可得到另一个 $C$ 中的点，
且与 $x$ 的距离不超过 $3^{-n}$。因此 $C$ 没有孤立点。零测集由此可以同时
是闭集、不可数集和完全集。

\section{正测度而无内点的集合}

把删去长度安排得更快收敛，可以保留正测度。Smith--Volterra 集的一个
构造从 $[0,1]$ 开始：第 $n\ge1$ 步从 $2^{n-1}$ 个剩余区间中各删去
长度 $4^{-n}$ 的居中开区间。记第 $n$ 步每个保留区间的长度为 $\ell_n$。
由 $\ell_0=1$ 和 $\ell_n=(\ell_{n-1}-4^{-n})/2$，归纳得到
\[
 \ell_n=2^{-n-1}+2^{-2n-1}.
\]
特别地，$\ell_{n-1}>4^{-n}$，每次删除都能在区间内部进行。
各步删去的区间两两不交，总删去长度为
\[
 \sum_{n=1}^{\infty}2^{n-1}4^{-n}=\frac12,
\]
故剩余闭集 $S$ 的测度为 $1/2$。若它包含非退化区间，该区间在每一阶段
都必须落在一个保留分支内，与 $\ell_n\to0$ 矛盾。因此 $S$ 无内点，且
$\partial S=S$。边界测度为正，所以上一章的判据表明 $S$ 不Jordan可测。
Cantor集与 $S$ 都闭且无内点，测度却分别为零与正数；仅靠这两个拓扑性质
不能判断集合是否零测。

\section{Vitali 不可测集}

不可测集的构造使用选择原则。在 $[0,1]$ 上定义等价关系
\[
 x\sim y\quad\Longleftrightarrow\quad x-y\in\Q.
\]
用选择公理从每个等价类选取一个代表，所得集合记为 $V$。对
$q\in\Q\cap[-1,1]$，令 $V_q=V+q$。

这些平移集两两不交：若 $v+q=w+r$，则 $v\sim w$，代表的唯一性给出
$v=w$，继而 $q=r$。另一方面，每个 $x\in[0,1]$ 与 $V$ 中某个代表等价，
所以
\[
 [0,1]\subseteq\bigcup_{q\in\Q\cap[-1,1]}V_q
 \subseteq[-1,2].
\]

\begin{theorem}[Vitali]
\label{thm:vitali-nonmeasurable-set}
上述代表集 $V$ 不是 Lebesgue 可测集。
\end{theorem}

\begin{proof}
若 $V$ 可测，平移不变性使所有 $V_q$ 具有同一测度 $c$。若 $c=0$，可数
次可加性给出 $[0,1]$ 的测度不超过零，矛盾。若 $c>0$，两两不交集族的并
测度为 $\sum_qc=+\infty$，却包含于有限测度区间 $[-1,2]$，仍矛盾。
\end{proof}

因此不可能在 $\mathcal P(\R)$ 上定义一个取非负扩充实值、同时满足下列三项的函数：延拓
区间长度、平移不变、对可数个不交集合可加。Carathéodory 构造给出一个
包含全部 Borel 集并且完备的 $\sigma$-代数；Vitali 定理说明，在保留上述
三项性质时，这个定义域不能扩大到全部子集。

Vitali 集的存在依赖从不可数多个等价类同时取代表。正文采用 ZFC 背景，
所以允许这一步；若研究较弱选择公理下的集合可测性，结论需要重新表述。

\section{不同的大小概念}

在 $\R^d$（$d\ge1$）中有蕴含链
\[
 \text{有限}\ \Longrightarrow\ \text{至多可数}\ \Longrightarrow\
 \text{Lebesgue零测}\ \Longrightarrow\ \text{无内点}.
\]
最后一步因为每个非空开球都含有正体积长方体。反向一般均不成立：
$\Q$、Cantor集和Smith--Volterra集依次给出反例。
以下表格的拓扑性质均相对于 $[0,1]$ 而言：
\[
\begin{array}{c|c|c|c}
\text{集合}&\text{基数}&\text{测度}&\text{拓扑性质}\\ \hline
\Q\cap[0,1]&\text{可数无限}&0&\text{稠密、无内点}\\
C&\text{不可数}&0&\text{闭、无内点}\\
S&\text{不可数}&1/2&\text{闭、无内点}\\
\{x\in[0,1]:x\notin\Q\}&\text{不可数}&1&\text{稠密、无内点}
\end{array}
\]

另一个常见的比较来自范畴。闭包无内点的集称为\term{无处稠密集}；
至多可数个无处稠密集的并称为\term{第一纲集}，其补称为\term{余第一纲集}。
$S$ 本身无处稠密而测度为正。相反，也有零测的余第一纲集。
枚举 $\Q=\{q_1,q_2,\ldots\}$，令
\[
 G_n=\bigcup_{j\ge1}(q_j-2^{-n-j-2},q_j+2^{-n-j-2}),
 \qquad G=\bigcap_{n\ge1}G_n.
\]
每个 $G_n$ 开且含 $\Q$，故稠密；$m(G_n)\le2^{-n-1}$，所以 $m(G)=0$。
补集 $\R\setminus G$ 是闭无处稠密集 $\R\setminus G_n$ 的可数并，
因而为第一纲集，且在每个有界区间中占满测度。

这里的 $G$ 并非第一纲集。为核对这一点，只需证明 $\R$ 不能是闭无处
稠密集 $F_n$ 的可数并。从一个非退化闭区间出发，逐步在前一区间内部
选取非退化闭子区间 $I_n$，使 $I_n\cap F_n=\varnothing$ 且长度至多
$2^{-n}$。无处稠密性保证选择可行，区间套交点避开全部 $F_n$。
对任意第一纲覆盖先将成员取闭包，就归结为这一情形。如果 $G$ 也是
第一纲集，它与已经为第一纲的补集将覆盖 $\R$，矛盾。
所以零测与第一纲这两种性质互不蕴含。

测度关注可数覆盖的总体积；拓扑关注邻域；基数只关心是否存在双射。以后
出现“几乎处处”时，只能删除零测集，不能把它替换为“可数多个例外”或
“没有内部的例外”。

\section*{习题}
\addcontentsline{toc}{section}{习题}

\subsection*{A组　零测与几乎处处}
\begin{enumerate}[label=\textbf{\thechapter.\arabic*.}]
 \exerciseitem{null-countable-union}
 证明可数个 Lebesgue 零测集的并仍为零测集。
 \exerciseitem{ae-equivalence-relation}
 证明函数之间“几乎处处相等”是等价关系。
 \exerciseitem{ae-equality-level-sets}
 设 $f,g$ 为 $\R^d$ 的同一Lebesgue可测子集 $X$ 上的有限实值或复值函数。
 证明 $f=g$ a.e. 当且仅当对每个 $n\ge1$，集合
 $\{|f-g|>1/n\}$ 零测。
 \exerciseitem{dense-null-set-construction}
 在任意非空开集 $O\subset\R^d$（$d\ge1$）中构造一个相对稠密的可数无限零测集。
\end{enumerate}

\subsection*{B组　Cantor 型集合}
\begin{enumerate}[label=\textbf{\thechapter.\arabic*.},resume]
 \exerciseitem{cantor-ternary-characterization}
 证明 $x\in C$ 当且仅当它有一个只含 $0,2$ 的三进制表示。
 \exerciseitem{cantor-perfect-nowhere-dense}
 证明 Cantor 集闭、无孤立点且无内点。
 \exerciseitem{cantor-cardinality}
 构造 $\{0,1\}^{\Npos}$ 到 Cantor 集的单射，并推出 $C$ 与 $[0,1]$
 等势。
 \exerciseitem{fat-cantor-measure}
 核对 Smith--Volterra 构造中各步区间互不重叠、总删去长度为 $1/2$，并
 证明剩余集无内点。
 \exerciseitem{closed-set-boundary-jordan}
 设 $F\subset[a,b]$ 闭且无内点。证明 $F$ Jordan 可测当且仅当
 $m(F)=0$。
\end{enumerate}

\subsection*{C组　不可测性}
\begin{enumerate}[label=\textbf{\thechapter.\arabic*.},resume]
 \exerciseitem{vitali-translates-disjoint}
 详细证明 Vitali 构造中的平移集 $V+q$ 两两不交及其覆盖包含关系。
 \exerciseitem{vitali-inner-outer-measure}
 对 $A\subseteq\R$ 定义内测度
 $m_*(A)=\sup\{m(E):E\subseteq A,\ E\text{为Lebesgue可测集}\}$。
 证明 Vitali 集的内测度为零，并说明其外测度不能为零。
 \exerciseitem{nonmeasurable-subset-null-impossible}
 说明完备性为何排除“零测集含不可测子集”，并由此解释 Vitali 集不可能
 包含在 Borel 零测集中。
\end{enumerate}
